\section{Raw-Gradient Diagnostics under Adam}
\label{sec:adam-diagnostics}

Adam~\cite{kingma2014adam} transforms the raw gradient through first- and
second-moment estimates,
\begin{equation}
\begin{aligned}
 m_t&=\beta_1m_{t-1}+(1-\beta_1)g_t,\\
 v_t&=\beta_2v_{t-1}+(1-\beta_2)g_t^{\odot2},
\end{aligned}
\label{eq:adam-moments}
\end{equation}
with zero initial moments. Bias correction gives
$\widehat m_t=m_t/(1-\beta_1^t)$ and
$\widehat v_t=v_t/(1-\beta_2^t)$, followed by
\begin{equation}
 d_t=\frac{\widehat m_t}{\sqrt{\widehat v_t}+\epsilon},
 \qquad \theta_t=\theta_{t-1}-\eta d_t,
 \label{eq:adam-direction}
\end{equation}
where operations are coordinatewise. Throughout this section, $g_t$ is
the raw loss gradient and $\rho_t$ remains the cosine between consecutive
raw gradients, as in Section~\ref{sec:displacement}.
We evaluate $\eta\widehat\kappa_t$ using the same diagnostic in
Eq.~\eqref{eq:kappa-estimate}. Although $d_t$ determines the step,
$g_t$ describes the local loss slope. Changes in its direction and norm
reveal how that slope varies across the states visited by Adam.
Comparing raw-gradient and first-moment measurements then shows how this
response propagates through the optimizer's memory. Momentum cosines
are labeled separately in the figures.

We reuse the CIFAR-10/MLP, CIFAR-10/ResNet-32 and WikiText-2/Transformer
objectives from Section~\ref{sec:gd-experiments}. Each Adam trajectory starts
from the same frozen, untrained seed-0 weights used for its GD comparison,
with fresh moments. We run 1,024 full-objective updates at each constant
rate: $\{0.0001,0.0003,0.001\}$ for the MLP and
$\{0.0003,0.001,0.003\}$ for ResNet and Transformer. All nine arms use
$(\beta_1,\beta_2)=(0.9,0.999)$ and $\epsilon=10^{-8}$, without a schedule,
clipping or weight decay. There is one initialization per model, matched
across rates. Full-batch accumulation removes sampling variation in the
objective; model, objective and numerical details are in
Appendix~\ref{appendix:adam-measurements}.

\begin{figure*}[tp]
 \centering
 \includegraphics[width=\textwidth]{figures/analysis036/adam_manuscript_overview.pdf}
 \caption{Full-batch Adam training from matched initial weights, with three
 constant learning rates per model. Every training-loss observation is shown
 without smoothing on logarithmic y-axes. Parameter state $s$ denotes $s$
 completed updates, from initialization to 1,024 updates. One initialization
 per model; colors identify increasing rates within each panel.}
 \label{fig:adam-overview}
\end{figure*}

Figure~\ref{fig:adam-overview} shows overall convergence with
transient loss spikes, especially at larger rates. The largest tested rate
attains the lowest final training loss in each model despite these excursions.
These excursions interrupt trajectories whose longer-term loss continues
to decrease. We next examine their gradient dynamics before summarizing
the spike distribution across the nine arms.

\subsection{Anatomy of Individual Spikes}

For each model, we select the recovered excursion with the largest
peak/pre-rise loss ratio at its highest tested rate. Selection uses only
loss, with no condition on a diagnostic. The selected peaks occur at
states 703, 422 and 794 for MLP, ResNet and Transformer, respectively.

\begin{figure*}[tp]
 \centering
 \includegraphics[width=\textwidth]{figures/analysis036/adam_manuscript_spikes.pdf}
 \caption{Largest recovered relative loss spike at each model's highest
 Adam rate. Rows show loss, raw-gradient projected change, consecutive
 cosines of raw gradients and first moments, and their relative norms.
 All observations are unsmoothed. At parameter state $s$, vector diagnostics
 refer to $g_{s+1}$ and $m_{s+1}$. Norms are divided by their value at the
 first displayed state. Gray bands mark the pronounced rise from the first
 +10\% crossing to the loss peak (697--703, 418--422 and 790--794), rather
 than the entire excursion through recovery. The horizontal dotted loss
 reference is the first displayed loss. Windows begin ten states before
 the peak and end at the first later loss below that displayed reference.}
 \label{fig:adam-spikes}
\end{figure*}

All three examples in Fig.~\ref{fig:adam-spikes} enter the pronounced rise
with negative raw-gradient cosine. During the rise, the gradient norm
increases sharply and $\eta\widehat\kappa$ exceeds two before the loss
maximum. Momentum norms also grow.

The descent separates the models. ResNet and Transformer return to
constructive regimes, with positive gradient and momentum cosines,
falling $\eta\widehat\kappa$ and a comparatively smooth loss decline.
The MLP continues to oscillate: raw cosine remains mostly negative,
momentum cosine alternates between signs, and loss and norms decrease
through repeated rebounds. Together, the examples expose developing
spikes through destructive alignment and sharp gradient growth.
The early increase in $\eta\widehat\kappa$ motivates intervention before
the loss reaches its maximum.

\subsection{Loss-Defined Events and Their Distribution}

To examine repetition without selecting on gradients, write
$\mathcal L_s:=\mathcal L(\theta_s)$ for the recorded training loss.
Starting at initialization, let $a_i$ be the earliest available state with
$\mathcal L_{a_i+1}>\mathcal L_{a_i}$. Hold the pre-rise reference
$\mathcal L_{a_i}$ fixed. Recovery occurs at the first later state whose
loss is at or below this pre-rise level, $\mathcal L_s\leq\mathcal L_{a_i}$.
For an event recovered within $T$ updates, define
\begin{equation}
\begin{aligned}
 c_i&:=\min\{s:a_i<s\leq T,\ \mathcal L_s\leq\mathcal L_{a_i}\},\\
 p_i&\in\operatorname*{arg\,max}_{a_i<s<c_i}\mathcal L_s.
\end{aligned}
\label{eq:adam-spike-boundaries}
\end{equation}
If several states attain the maximum, $p_i$ is the earliest one.
Figure~\ref{fig:adam-spike-definition} illustrates these boundaries.
An event is included only when recovery is observed by $T=1024$.
Search resumes at $c_i$; partial declines and additional local maxima
before this return belong to the same event. Thus $a_i$ is the last
pre-rise state, $a_i+1$ the first elevated state, $p_i$ the loss peak,
and $c_i$ the first recovered state. Peak increase is
$100(\mathcal L_{p_i}/\mathcal L_{a_i}-1)$; rise and recovery times are
$p_i-a_i$ and $c_i-p_i$, in optimizer updates. The three loss ratios are
$1$, $\mathcal L_{p_i}/\mathcal L_{a_i}$ and
$\mathcal L_{c_i}/\mathcal L_{a_i}\leq1$.

\begin{figure*}[tp]
 \centering
 \begin{minipage}[c]{0.49\textwidth}
 \includegraphics[width=\columnwidth]{figures/analysis036/adam_spike_definition.pdf}
 \end{minipage}\hfill
 \begin{minipage}[c]{0.47\textwidth}
 \caption{Spike boundaries (schematic). The pre-rise state $a_i$ fixes
 the reference loss; $a_i+1$ is the first elevated state, $p_i$ the peak,
 and $c_i$ the first return to or below that reference. Rise and recovery
 durations are $p_i-a_i$ and $c_i-p_i$.}
 \label{fig:adam-spike-definition}
 \end{minipage}
\end{figure*}

\begin{figure*}[tp]
 \centering
 \includegraphics[width=\textwidth]{figures/analysis036/adam_peak_increase_density.pdf}
 \caption{Distributions of peak percentage increases for all 806 recovered
 excursions, including small fluctuations, with no amplitude cutoff.
 Colors identify LRs within each model. Gaussian kernel densities are
 estimated in $\log_{10}$ percentage-increase coordinates with one pooled
 bandwidth per model, common across its LRs. Each LR density integrates
 to one in that coordinate. The log x-axis shows percentage increase:
 1 denotes +1\%, and 100 denotes +100\%. This differs from the loss-ratio
 axis in Fig.~\ref{fig:adam-rho-norm-spikes}.}
 \label{fig:adam-spike-distribution}
\end{figure*}

Figure~\ref{fig:adam-spike-distribution} places the illustrations within
the complete distribution of recovered excursions. Small fluctuations
coexist with long tails; larger rates populate more severe excursions,
although the distributions overlap and do not undergo a uniform shift.
To focus the anatomy analysis on pronounced changes, we retain events
with peak increase strictly greater than 10\%.
The filtered sample contains 129 recovered events. Two further
above-threshold candidates without observed recovery are excluded.

\begin{table*}[tp]
 \caption{Anatomy of recovered spikes with peak increase greater than 10\%.
 Cells report mean $\pm$ sample SD across events; maximum peak increase has
 its own subcolumn. Before denotes a single diagnostic at state $a_i$;
 recovery denotes the event's mean over states $p_i+1,\ldots,c_i$.
 Norm ratios compare loss-peak/pre-rise and recovered/loss-peak endpoints.
 $N$ counts included events; dashes indicate no qualifying event.
 Two initialization events have undefined pre-rise diagnostics, leaving
 46 and 38 observations for those entries in the MLP's middle and highest
 LR rows. No smoothing or averaging across initializations.}
 \label{tab:adam-spike-anatomy}
 \centering
 \footnotesize
 \setlength{\tabcolsep}{1.5pt}
 \renewcommand{\arraystretch}{1.15}
 \input{tables/analysis036_spike_anatomy}
\end{table*}

\begin{figure*}[tp]
 \centering
 \includegraphics[width=\textwidth]{figures/analysis036/adam_rho_gradient_rise_peak_scatter.pdf}
 \caption{Pre-rise raw-gradient cosine times rise norm ratio versus
 peak/pre-rise loss for 127 recovered spikes above 10\%; two initialization
 events have undefined pre-rise cosine. One unsmoothed point per event.
 Colors identify LRs. Separate per-LR least-squares fits with intercept use
 $\log_{10}$(peak/pre-rise loss) as the response; color-matched $R^2$ values
 are descriptive in-sample scores on that log response. Dashed lines extend
 across each panel, including extrapolation beyond observed samples.
 The y-axis is a loss ratio: 1 denotes no increase and 2 denotes +100\%.}
 \label{fig:adam-rho-norm-spikes}
\end{figure*}

For Table~\ref{tab:adam-spike-anatomy}, pre-rise diagnostics are measured
at $\theta_{a_i}$, hence use gradient index $a_i+1$. Recovery entries
first average each event's diagnostic over the arrival states
$p_i+1,\ldots,c_i$, then summarize those means across events with equal
event weight. They are not values only at $c_i$. The norm-ratio columns
instead use the endpoints
$\|g_{p_i+1}\|/\|g_{a_i+1}\|$ and
$\|g_{c_i+1}\|/\|g_{p_i+1}\|$. The gradient norm's own maximum need not
coincide with the loss peak.

Among qualifying events, both the mean increase and the maximum increase
grow with LR in every model with more than one eligible rate. The largest
increase is 858.8\% for ResNet at $\eta=0.003$. Mean recovery takes
17--19 updates versus 6--8 for the ResNet rise, and 24--31 versus about
6 for the Transformer. MLP rise and recovery times are more comparable.
Thus, a short amplification can be followed by a much longer descent,
without implying that all intermediate loss changes are monotonic.

Pre-rise raw cosine is negative in 92.1\% of events with a defined
pre-rise measurement. The gradient norm at the loss peak exceeds its
pre-rise value in 100\% of events. Projected change is already above two
before the first loss increase in 16.5\% of defined cases. Its full-event
maximum precedes the loss peak in 63.6\% of events, coincides in 22.5\%
and follows in 14.0\%.

\subsection{Reversal, Norm Amplification and Possible Interventions}

To relate the initial regime to amplification during an event, we plot
\begin{equation}
 \rho_{a_i+1}\frac{\|g_{p_i+1}\|}{\|g_{a_i+1}\|}
 \quad\text{against}\quad
 \frac{\mathcal L_{p_i}}{\mathcal L_{a_i}}.
 \label{eq:adam-event-product}
\end{equation}
This combines the pre-rise consecutive cosine with the subsequent rise
norm ratio. Figure~\ref{fig:adam-rho-norm-spikes} shows that a destructive
initial regime combined with gradient amplification explains a substantial
part of spike magnitude: the per-rate fits account for 67--69\% of
log-loss-ratio variation for the MLP, 90--94\% for ResNet, and 38--54\%
for the Transformer. This relationship connects the individual spike
examples to the broader event sample. Reversal identifies the regime;
gradient growth explains how strongly that reversal is amplified.

This anatomy motivates testing whether intervention during reversal and
early norm growth can attenuate the subsequent excursion. Gradient-norm
clipping was proposed to address exploding gradients~\cite{pascanu2013difficulty},
and analyses allowing smoothness to grow with gradient norm provide a
further rationale for adaptive scaling~\cite{zhang2020clipping}. Learning-rate
adjustment and data-mixing interventions provide other tests of the proposed
interaction between the current regime and subsequent gradient growth.
These tests should isolate states before pronounced excursions, compare
matched interventions with unchanged continuations, and measure both spike
attenuation and later loss at equal budgets. Establishing prospective
warning information also requires non-spike comparisons. Such interventions
and predictive validation remain future work.
